%% GENERATED FILE -- do not hand-edit.
%% SOURCE: experiments/019-perturb-seq-costing/scripts/render_tex_tables.py
\begingroup\scriptsize
\setlength{\LTleft}{0pt}\setlength{\LTright}{0pt}
\begingroup\captionsetup{type=table}
\captionof{table}[]{Provenance for every number drawn by hand in Figs.~\ref{fig:methods-map}--\ref{fig:droplet}. The matplotlib figures are omitted because their generating scripts read the data directly. A \textbf{note} marks a number that needed a judgement call or that was corrected in review; an entry with no citation key is not backed by the library mirror and should not be quoted. Generated from \file{experiments/019-perturb-seq-costing/scripts/figure_sources.py}.}\label{tab:figure-provenance}
\endgroup
\begin{longtable}{@{}>{\raggedright\arraybackslash}p{0.20\textwidth}>{\raggedright\arraybackslash}p{0.15\textwidth}>{\raggedright\arraybackslash}p{0.60\textwidth}@{}}
\toprule
Element & As drawn & Source \\
\midrule
\endfirsthead
\multicolumn{3}{@{}l}{\scriptsize\itshape Table~\ref{tab:figure-provenance}, continued}\\
\toprule
Element & As drawn & Source \\
\midrule
\endhead
\bottomrule
\endlastfoot
\multicolumn{3}{@{}l}{\bfseries scrnaseq-method-taxonomy.drawio} \\*[1pt]
Microfluidic chip (C1) throughput & \texttt{up to 96 per chip} & ``It uses an integrated fluidic circuit to isolate and process up to 96 cells at once.'' \citep{nadal-ribellesRiseSinglecellTranscriptomics2024} {\color{tcgray}\emph{Note:} Was drawn as \textasciitilde{}160 cells, which is Gasch et al.'s 163-cell study total across four chip runs -- a study total on a per-experiment axis. The capacity of one chip is the comparable quantity.} \\
\addlinespace[2.5pt]
FACS into plates throughput & \texttt{\textasciitilde{}285 cells} & ``we applied yscRNA-seq to 285 individual yeast cells'' \citep{nadal-ribellesSensitiveHighthroughputSinglecell2019} \\
\addlinespace[2.5pt]
Microdissection throughput & \texttt{\textasciitilde{}2,000 cells} & ``Microdissection coupled with imaging has been used to profile the transcriptomes of Schizosaccharomyces pombe across a variety of environmental stressors and S. cerevisiae during aging (Saint et al., 2019; Wang et al., 2022)'' \tcflagext {\color{tcgray}\emph{Note:} NOT SOURCED TO A NUMBER. The review names the studies but states no cell count, and neither primary is in the mirror. The \textasciitilde{}2,000 is an order-of-magnitude placeholder and is the one number in this figure that should not be quoted.} \\
\addlinespace[2.5pt]
Droplet throughput range & \texttt{6,000-100,000} & ``applied it to over 100,000 single cells from three crosses [upper bound; lower bound is Jariani et al.'s 6,118 cells]'' \citep{boocockSinglecellEQTLMapping2025} \\
\addlinespace[2.5pt]
Microwell array throughput & \texttt{up to 1,061,865} & ``we used a microwell-based platform for single-cell isolation ... We profiled a total of 1.061.865 cells'' \citep{nadal-ribellesSinglecellResolvedGenotypephenotype2025} {\color{tcgray}\emph{Note:} Was \textasciitilde{}75,000, taken from mDrop-seq -- which is a droplet method, not a microwell one. Replaced with the microwell study that exists.} \\
\addlinespace[2.5pt]
Combinatorial indexing throughput & \texttt{25,000-240,000 per run} & ``Reverse transcription (in situ) was performed in 25 uL reactions in 48 wells of a 100 uL 96-well plate ... 200,000 cells/mL [= \textasciitilde{}5,000 cells x 48 wells = \textasciitilde{}240,000 barcoded per run; lower bound is Kuchina et al.'s 25,214-cell combined dataset]'' \citep{brettnerUltraHighthroughputMassively2024} {\color{tcgray}\emph{Note:} The upper bound was 1,000,000, which is Gaisser et al.'s protocol capacity claim ("enables ... up to 1 million"), not a profiled count. Same correction as Fig. 4.} \\
\addlinespace[2.5pt]
\multicolumn{3}{@{}l}{\bfseries splitseq-barcoding.drawio} \\*[1pt]
round-1 wells and cells per well & \texttt{48 wells, \textasciitilde{}5,000 cells each} & ``Reverse transcription (in situ) was performed in 25 uL reactions in 48 wells of a 100 uL 96-well plate with each well containing wellspecific, barcoded primers ... 200,000 cells/mL'' \citep{brettnerUltraHighthroughputMassively2024} \\
\addlinespace[2.5pt]
barcode space over three rounds & \texttt{963 = 884,736} & ``Cells are then pooled and randomly split into a new 96-well plate, and a well-specific barcode is attached'' \citep{brettnerUltraHighthroughputMassively2024} {\color{tcgray}\emph{Note:} Protocol capacity. The published run loaded 48 wells in round 1, so its realized space was 48x96x96 = 442,368 (Sec. 3.1).} \\
\addlinespace[2.5pt]
sublibrary size & \texttt{5,000-20,000 cells} & ``In experiment 3, the barcoded cells were evenly divided into 10 sublibraries. The two sequenced sublibraries returned \textasciitilde{}5,500 and 10,000 barcoded cells that passed computational filtering.'' \citep{brettnerUltraHighthroughputMassively2024} \\
\addlinespace[2.5pt]
rRNA fraction of recovered transcripts & \texttt{93.7\%} & ``On average, the transcript proportions we recover are 93.7\% rRNA, 0.005\% tRNA, 0.04\% ncRNA, and 5.75\% mRNA.'' \citep{brettnerUltraHighthroughputMassively2024}\,\texttt{(md:70)} \\
\addlinespace[2.5pt]
\multicolumn{3}{@{}l}{\bfseries droplet-10x-barcoding.drawio} \\*[1pt]
doublet rate at 20,000 cells & \texttt{\textasciitilde{}8\% at 20,000 recovered} & ``[10x Chromium specification for the GEM-X/v3.1 chemistry priced in Sec. 5; the protocol paper used 3' v2]'' \citep{jarianiNewProtocolSinglecell2020} {\color{tcgray}\emph{Note:} MANUFACTURER SPECIFICATION, not a measurement from a mirrored paper. Verify against the current 10x user guide before quoting.} \\
\addlinespace[2.5pt]
digestion timing and cell-count drop & \texttt{6 min hold, 200-fold drop within 20 min at 53 C} & ``cell counts hold on ice and through the 6 min of droplet generation, then drop 200-fold within 20 min at 53 C'' \citep{jarianiNewProtocolSinglecell2020} \\
\addlinespace[2.5pt]
bead barcode and UMI lengths & \texttt{16 nt CB, 12 nt UMI, 118 cycles} & ``[10x 3' v3/GEM-X read layout; cycle count derived in experiments/019-perturb-seq-costing/scripts/read\_structure.py]'' \citep{jarianiNewProtocolSinglecell2020} {\color{tcgray}\emph{Note:} Layout is the current GEM-X/v3.1 chemistry, not the v2 Jariani ran.} \\
\addlinespace[2.5pt]
\end{longtable}
\endgroup
